\documentclass[../differential_methods.tex]{subfiles}
\begin{document}
\section{Aula 04 - 11 de Março, 2026}
\subsection{Motivações}
\begin{itemize}
	\item Equações não lineares;
	\item Discretização de Problemas não Lineares;
	\item Método de Newton.
\end{itemize}
\subsection{Equações não lineares}
O primeiro exemplo de uma equação não linear, e um dos mais famosos, é o movimento de um pêndulo: considere o movimento de um pêndulo com uma partícula de massa m presa ao final de uma haste de comprimento L, considerada aqui sem massa, e seja \(\theta (t)\) o ângulo do pêndulo verticalmente no momento t.
Ignorando a massa da barra, forças de atrito e a resistência do ar, o movimento deste pêndulo pode ser descrito como
\[
	\theta ''(t)=-\frac{g}{L}\sin^{}{(\theta (t))},
\]
sendo g a constante gravitacional. Aplicando a segunda lei de Newton às forças atuando na partícula presa à ponta, obtemos a relação
\[
	F = m a = -mg\sin^{}{(\theta )} \Rightarrow a = -g \sin^{}{(\theta )}.
\]
Como o comprimento de arco traçado quando o ângulo vale \(\theta \) é dado por \(s = L \theta \), observamos também que
\[
	v = \frac{\mathrm{d}s}{\mathrm{d}t} = L \frac{\mathrm{d}\theta }{\mathrm{d}t},
\]
e consequentemente
\[
	a = \frac{\mathrm{d}^{2}s}{\mathrm{d}t^{2}} = L \frac{\mathrm{d}^{2}\theta }{\mathrm{d}t^{2}}.
\]
Portanto, a descrição analítica do problema é que
\[
	\frac{\mathrm{d}^{2}\theta }{\mathrm{d}t^{2}}=-\frac{g}{L}\sin^{}{(\theta )}.
\]

Nosso objetivo é encontrar uma forma da solução desse sistema na qual possamos testar nossos métodos numéricos, porque o caso aqui vai ser bem diferente do que vimos anteriormente. Sendo assim, começando pela resolução do sistema, vamos supor que \(\frac{g}{L}=1\) para simplificar; com isso, temos
\begin{align*}
	 & \theta ''(t)=-\sin^{}{(\theta (t))}, \quad 0<t<T \\
	 & \theta (0)=\alpha                                \\
	 & \theta (T)=\beta ,
\end{align*}
mas vale observar que fornecer um valor da velocidade angular \(\theta'(0)\) é mais natural em contextos reais. Para amplitudes pequenas do ângulo \(\theta \), é costume utilizar a aproximação \(\sin^{}{(\theta (t))}\approx \theta (t)\), tornando o problema acima em
\begin{align*}
	 & \theta ''(t)=\theta (t), \quad 0<t<T \\
	 & \theta (0)=\alpha                    \\
	 & \theta (T)=\beta ,
\end{align*}
que sabemos resolver com soluções gerais da forma
\[
	\theta (t)=A\cos^{}{(t)}+B\sin^{}{(t)}.
\]

Seguindo a abordagem dos problemas lineares, podemos aproximar os valores das soluções por
\[
	\frac{1}{h^{2}}(\theta_{j-1}-2\theta_{j}+\theta_{j+1})+\sin^{}{(\theta_{j})}=0, \quad j=1,2,\dotsc ,m,
\]
onde
\[
	h\coloneqq \frac{T}{(m+1)},\quad \theta_{0}\coloneqq \alpha  \quad\&\quad \theta_{m+1}\coloneqq \beta .
\]
Transformamos o problema em um sistema de equações não lineares da forma
\[
	G(\theta )=0,
\]
onde \(G:\mathbb{R}^{m}\rightarrow \mathbb{R}^{m}\) e \(\theta = (\theta_1, \theta_2, \dotsc , \theta_{m})^{T}\),
que tem um problema grande de não poder ser resolvido utilizando sistemas lineares tridiagonais: ao invés de usarmos métodos diretos, é necessário explorarmos métodos iterativos, tal qual o \textit{Método de Newton}.

A forma como o método de Newton funciona é que, se \(\theta^{[k]}\) denota a aproximação para \(\theta \) no k-ésimo passo, então o método de Newton consiste em utilizar a expansão de Taylor
\[
	G(\theta^{[k+1]})=G(\theta^{[k]})+\frac{\partial^{}}{\partial \theta ^{}}G(\theta^{[k]})(\theta^{[k+1]}-\theta^{[k]})+\dotsc,
\]
colocar \(G(\theta^{[k+1]})=0\) como desejado, e abandonar os termos de ordens maiores, resultando em
\[
	0\approx G(\theta^{[k]})+\frac{\partial^{}}{\partial \theta ^{}}G(\theta^{[k]})(\theta^{[k+1]}-\theta^{[k]}).
\]
Disso, podemos fazer uma etapa do método conhecida como atualização, colocando
\[
	\theta^{[k+1]}\coloneqq \theta^{[k]}+\delta^{[k]},
\]
onde \(\delta^{[k]}\) é a solução da equação com Jacobiana
\[
	J(\theta^{[k]})\delta^{[k]}=-G(\theta^{[k]} ),
\]
onde
\[
	J(\theta )\coloneqq \frac{\partial^{}}{\partial \theta ^{}}G(\theta )\in \mathbb{R}^{m\times m}, \quad J_{ij}(\theta )=\frac{\partial^{}}{\partial \theta_{j}^{}}G_{i}(\theta ),\; i, j=1,2,\dotsc ,m.
\]

Em nosso exemplo do pêndulo,
\[
	G_{i}(\theta )=G_{i}(\theta_1, \theta_2, \dotsc, \theta_{m})=\frac{1}{h^{2}}(\theta_{i-1}-2\theta_{i}+\theta_{i+1})+\sin^{}{(\theta_{i})},
\]
tal que
\[
	J_{ij}(\theta )  = \left\{\begin{array}{ll}
		\frac{1}{h^{2}},                        & \quad j=i-1 \text{ ou } j=i+1, \\
		-\frac{2}{h^{2}}+\cos^{}{(\theta_{i})}, & \quad j=1,                     \\
		0,                                      & \quad \text{caso contrário}
	\end{array}\right.,
\]
resultando na matriz
\[
	A=\frac{1}{h^{2}} \begin{bmatrix}
		-\frac{2}{h^{2}}+\cos^{}{(\theta_1)} & 1                                      & \dotsc &        &                                        \\
		1                                    & -\frac{2}{h^{2}}+\cos^{}{(\theta_{2})} & 1      & \dotsc &                                        \\
		\vdots                               & \ddots                                 & \ddots & \cdots &                                        \\
		                                     &                                        &        & 1      & -\frac{2}{h^{2}}+\cos^{}{(\theta_{m})}
	\end{bmatrix}.
\]

\begin{tcolorbox}[
		skin=enhanced,
		title=Observação,
		fonttitle=\bfseries,
		colframe=black,
		colbacktitle=cyan!75!white,
		colback=cyan!15,
		colbacklower=black,
		coltitle=black,
		drop fuzzy shadow,
		%drop large lifted siador
	]
	O método de Newton requer uma ressalva: o chute inicial deve ser perto o suficiente do valor real, e ele converge de forma quadrática se isso ocorrer!
\end{tcolorbox}
A solução do problema não linear acima é uma \textit{solução isolada}, no sentido que não há outras soluções em uma vizinhança dela: ela é \textit{localmente única}, mas essa \textit{não} é a única solução ao problema de valor de contorno.

É importante manter em mente que há uma diferença entre a convergência do Método de Newton aplicado à solução de uma equação de diferença finitas e da convergência da aproximação de diferença finita da solução: inserindo a solução verdadeira à equação de diferenças finitas,
\begin{align*}
	\tau_{i} & \coloneqq \frac{1}{h^{2}}(\theta(t_{i-1})-2\theta(t_{i}))+\sin^{}{(\theta(t_{i}))}-0                                           \\
	         & = \biggl[\theta''(t_{i})+\sin^{}{\biggl(\theta(t_{i})+\frac{1}{12}h^{2}\theta^{(4)}(t_{i})\biggr)}\biggr] + \mathcal{O}(h^{4}) \\
	         & = \frac{1}{12}h^{2}\theta^{(4)}(t_{i})+\mathcal{O}(h^{4}),\quad i=1,2,\dotsc ,m.
\end{align*}
Concluímos, então, que o erro de truncamento local para a aproximação é da ordem de \(h^{2}\).

Seja \(\hat{\theta }\) o vetor de valores reais numa malha de pontos, e \(\tau=(\tau_1, \tau_2, \dotsc , \tau_{m})^{T}\). Então, se \(G(\theta )=\tau \) e \(E\coloneqq \theta -\hat{\theta }\) é o erro global, temos
\[
	G(\theta )-G(\hat{\theta })=0-\tau = -\tau.
\]
No caso linear, onde \(G(\theta )=A\theta - F\), segue que
\[
	G(\theta )-G(\hat{\theta })=A\theta - A\hat{\theta } = -\tau \Rightarrow \dotsc \Rightarrow E^{h} = (A^{h})^{-1}(-\tau^{h});
\]
com uma expansão de séries de Taylor, podemos escrever
\[
	G(\theta )=G(\hat{\theta })+J(\hat{\theta })E + \mathcal{O}(\Vert E \Vert^{2}) \Rightarrow J(\hat{\theta })E = -\tau + \mathcal{O}(\Vert E \Vert^{2}),
\]
e ignorando os termos de ordem maior, então novamente temos um relação linear entre os erros locais e globais! Adotaremos a notação
\[
	\hat{J^{h}}\coloneqq J(\hat{\theta })
\]
em uma malha com espaçamento h.

\begin{def*}
	O método de diferenças não lineares \(G(\theta )=0\) é \textbf{estável em uma norma} \(\Vert \cdot  \Vert\) se as matrizes \((\hat{J^{h}})\) forem uniformemente limitadas nela, conforme h tende a zero; noutras palavras, se existe \(C>0\) e \(h_{0}>0\) tais que
	\[
		\biggl\Vert (\hat{J^{h}})^{-1} \biggr\Vert \leq C, \quad 0<h<h_{0}. \; \square
	\]
\end{def*}
Pode-se mostrar que, se o método é estável nesse sentido e consistente no sentido de que \(\Vert \tau^{h} \Vert\) tende a zero conforme h tende também, então o método converge (\(\Vert E^{h} \Vert\to 0\) conforme \(h\to0^{+}\)).

\end{document}
